\section{The remaining three-row generating functions}\label{app:threegf}

This appendix completes the proof of Theorem~\ref{thm:threegf}, with the notation of
Section~\ref{sec:three} and $D=LN=1-2z-4z^3-4z^4$. A \emph{restricted run} is a string of
singletons with no two letters $b$ at distance two. By reversal and reflection the next lemma covers the classes $BA$, $BC$ and $CB$ too.

\begin{lemma}[forbidden windows in the classes containing $B$]\label{lem:threewindows}
Let $S$ be a candidate in the ordered class $AB$ or in the triple class, and let $s$ be an
arbitrary singleton. Then a window fails if and only if it is listed for its class below.
\begin{center}\small
\begin{tabular}{@{}lll@{}}
\toprule
class & windows & forbidden \\
\midrule
$AB$, $X$ & singletons only & $bab$, $bbb$, $bcb$ \\
$AB$ & one double & $bsA$, $Asb$, $bAb$, $bAc$, $cAb$, $aBa$, $aBb$, $bBa$, $cBc$ \\
 & two doubles & $bAB$, $ABa$ \\
$X$ & one triple & $bsX$, $Xsb$, $aXa$, $aXb$, $bXa$, $bXc$, $cXb$, $cXc$ \\
\bottomrule
\end{tabular}
\end{center}
\end{lemma}

\begin{proof}
The tests (a) to (e) are those of Lemma~\ref{lem:diamond}, and
Lemma~\ref{lem:dichotomy}(b) gives the singleton windows.

\emph{A double or triple in the middle.} A window $sAt$ behaves as in
Lemma~\ref{lem:sixofnine}(a), except that test (e) fails $bAb$ through $1\in B_{i+1}$ and
$B\in F$. In $sBt$ and $sXt$ column $i+1$ is black. The diamond is empty if $s=t=b$. If
$s=b\ne t$ it shows $(i+2,2)$ alone and fails by test (e) exactly when $\{2,t\}\in F$, and
symmetrically for $t=b\ne s$. If $s$ and $t$ are outer it shows the horizontal pair and fails
by test (c) exactly when $s=t$.

\emph{Two doubles.} In $uAB$ with $u=b$ the diamond shows $(i+2,2)$ and $(i+1,3)$, and test (d)
fails it since $(i+2,3)$ and $(i+1,2)$ are black. For $u\ne b$ three corners show. In $ABw$ it
is empty for $w=b$ and otherwise fails by test (e) exactly for $w=a$.

\emph{A double or triple at a tip.} In $Asb$ and $Xsb$ the diamond shows the vertical pair or
one point of column $i+1$, and tests (b) and (e) fail it since $B\in F$. Every other window
with a double or $X$ at the left tip passes. It shows three corners or an oblique pair whose
test (d) needs the white point $(i+1,2)$. Reversal gives the windows with the letter at the
right tip.
\end{proof}

We count restricted runs with a transfer on the last two bits. Give an outer letter the bit
$0$ and the weight $2$, and $b$ the bit $1$. Let $h_{ij}$ count,
by $z^n$, the continuations of length $n\ge0$ of a restricted run ending in the bits $i,j$. A letter $b$ is forbidden after the bits
$1,j$. So $h_{00}=1+z(2h_{00}+h_{01})$, $h_{01}=1+z(2h_{10}+h_{11})$, $h_{10}=1+2zh_{00}$ and
$h_{11}=1+2zh_{10}$, and eliminating gives $Dh_{00}=1+z+3z^2+2z^3$.

Read runs $w_1\cdots w_n$ with
$n\ge1$ outward from a double. Let $g_0$ count those with $w_1$ a prescribed outer letter and
$g_1$ those with $w_1=b$. Let $u$ and $v$ count the same runs with $w_2$, if present, outer. A fictitious outer letter before $w_1$ restricts nothing, so $g_0=zh_{00}$,
$g_1=zh_{01}$, $u=z+2z^2h_{00}=zh_{10}$ and $v=z+2z^2h_{10}=zh_{11}$. Hence
\begin{equation}\label{eq:runs}
\begin{aligned}
g_0&=\frac{z(1+z+3z^2+2z^3)}D,& g_1&=\frac{z(1+z)}D,\\
u&=\frac{z(1+2z^2+2z^3)}D,& v&=\frac zD.
\end{aligned}
\end{equation}

\begin{proposition}[the triple class]\label{prop:GX}
With~\eqref{eq:runs}, the function $G_X(z)=z\bigl(1+4u+2v+2u^2+v^2\bigr)$ has a denominator
dividing $D^2$.
\end{proposition}

\begin{proof}
In $UXW$ the runs are restricted, and $bsX$ and $Xsb$ make their second letters, read from
$X$, outer (Lemma~\ref{lem:threewindows}). So one nonempty run gives $2u+v$ on either side. With two, the window at $X$ allows only the neighbor pairs $(a,c)$, $(b,b)$ and
$(c,a)$, which give $2u^2+v^2$.
\end{proof}

A word of $AB$ is $UAVBW$ with restricted runs $U$, $V$ and $W$. By
Lemma~\ref{lem:threewindows}, a nonempty $U$ read from $A$ has an outer second letter, so it is
counted by $u$ or $v$. A nonempty $W$ read from $B$ has no such condition, so it is counted by
$g_0$ or $g_1$. Let $\ell_t$ count the allowed $U$, the empty one included, when a singleton $t$
stands right of $A$. Let $r_t$ count the allowed $W$ when $t$ stands left of $B$. The windows
$bAb$, $bAc$, $cAb$, $aBa$, $aBb$, $bBa$ and $cBc$ give
\begin{equation}\label{eq:lr}
\begin{aligned}
(\ell_a,\ell_b,\ell_c)&=(1+2u+v,\ 1+u,\ 1+2u),\\
(r_a,r_b,r_c)&=(1+g_0,\ 1+g_0+g_1,\ 1+g_0+g_1).
\end{aligned}
\end{equation}
For $s\in\{00,10\}$ and a bit $\beta$ let $J_{s\beta}$ count, as for $h$, the continuations
from $s$ whose last letter has the bit $\beta$. Solving as for $h$ gives
\begin{equation}\label{eq:J}
J=\frac1D\begin{pmatrix}1+2z^2+2z^3&z+z^2\\1&2z^2+2z^3\end{pmatrix}.
\end{equation}

\begin{proposition}[the class $\{A,B\}$, polynomial form computer-assisted]\label{prop:GAB}
With~\eqref{eq:runs}, \eqref{eq:lr} and~\eqref{eq:J},
\begin{equation}\label{eq:GAB}
\begin{aligned}
G_{AB}(z)={}&z^2(1+2u)(1+g_0+g_1)+z^3(\ell_ar_a+\ell_br_b+\ell_cr_c)\\
&+z^4\,(\ell_a+\ell_c,\ \ell_b)\,J\begin{pmatrix}r_a+r_c\\2r_b\end{pmatrix}.
\end{aligned}
\end{equation}
Moreover $G_{AB}=P_B/D^3$ with
\begin{equation}\label{eq:PB}
P_B(z)=z^2(1+z)(1+z^2-3z^3+4z^6+4z^7),
\end{equation}
and $P_B\bmod L=\tfrac z2$ and $P_B\bmod N=\tfrac14-\tfrac z2$.
\end{proposition}

\begin{proof}
The runs cover the windows inside them, and we split by the length $k$ of $V$. For $k=0$ the windows $uAB$ and $ABw$ require $u\ne b$ and
$w\ne a$, which gives the first term. For $k=1$ the window $AtB$ passes, and $uAt$ and $tBw$
give $\ell_t$ and $r_t$. For $k\ge2$ the window $Av_1v_2$ makes $v_2$ outer, $v_{k-1}v_kB$ passes,
and the weight is $\ell_{v_1}r_{v_k}$. After $v_2$ the bits are $00$ or $10$ as $v_1$ is outer or
$b$, and the rest is a continuation, in which $v_2$ keeps its factor $2$ for $k\ge3$. Summing
$r_{v_k}$ over the colors of the last letter gives $J_{s0}(r_a+r_c)+J_{s1}\,2r_b$, with $k=2$ as
the constant term of $J$.

Clearing denominators is computer-assisted. The script \nolinkurl{three_row_closed_form.py}
rebuilds $G_{AB}$ in SymPy from~\eqref{eq:runs}, \eqref{eq:lr} and~\eqref{eq:J} and
finds~\eqref{eq:PB}. The first coefficients $7$, $31$, $120$, $449$ and $1622$ agree with the
class counts from the definition (\texttt{three\_row\_definition\_check.json}). The remainders
follow from

\begin{align*}
P_B&=L\bigl(-2z^8-2z^7-z^6+z^4-\tfrac z2\bigr)+\tfrac z2,\\
P_B&=N\bigl(2z^8+4z^7+z^6-2z^5-2z^4+\tfrac32z^2+\tfrac z2-\tfrac14\bigr)+\tfrac14-\tfrac z2.
\end{align*}
\end{proof}

\begin{theorem}[the exact denominator, computer-assisted]\label{thm:denominator}
The reduced denominator of $\sum_{m\ge3}a_3(m)z^m$ is $Q=R^3L^3N^3$, with five triple poles.
\end{theorem}

\begin{proof}
By Theorem~\ref{thm:threegf}(a) the denominator divides $R^3D^3=Q$. The factors are coprime,
since $L-N=-2z-4z^2$ vanishes only at $0$ and $-\frac12$, where $N\ne0$, and $L(\frac13)=\frac19$,
$N(\frac13)=\frac{11}9$. The numerator in~\eqref{eq:GAC} is $\frac4{81}$ at $z=\frac13$, where
$G_{AB}$ and $G_X$ are analytic. At a root $\zeta$ of $L$, $P_B(\zeta)=\zeta/2\ne0$, so $G_{AB}$
has a pole of order three, $G_X$ one of order at most two, and $G_{AC}$ none. At a root of $N$,
$\zeta^2=-\frac12$ and $P_B(\zeta)=\frac14-\frac\zeta2\ne0$, and the same argument applies.
\end{proof}
